\section{Additional Proofs}

\subsection{Proof of Proposition~\ref{prop:linear=cos}}\label{app:proof:prop:linear=cos}


\begin{proof}
Let $n=\med(m, k; \fcos)$ and there exists $X=\{x_1,\dots,x_m\}\subset\R^{n}\setminus\{0\}$ \footnote{If there is $x_i = 0$, the proof is basically the same, we omit this situation for convenience.}that is $k$-shattered by $\fcos$.
Define the radial projection $\rho:\R^{n}\setminus\{0\}\to S^{n-1}$ by $\rho(x)=\frac{x}{\|x\|_2}$, and let
$Y=\{\rho(x_i)\}_{i=1}^m\subset S^{n-1}$.

By $k$-shattering, for each $S\subseteq X$ with $|S|\le k$ there exist $\vw_S\in\R^{n}\setminus\{0\}$ such that
\[
x\in S \Rightarrow r_S-\frac{\langle \vw_S,x\rangle}{\|\vw_S\|\|x\|} \le 0,
\qquad
x\in X\setminus S \Rightarrow r_S-\frac{\langle \vw_S,x\rangle}{\|\vw_S\|\|x\|} > 0.
\]

For $z\in S^{n-1}$ define the affine functional
\[
f_S(z)\;:=\;\big\langle \tfrac{-\vw_S}{\|\vw_S\|},\, z \big\rangle + r_S
\ \in\ \fln.
\]

Then, for every $x\in X$ we have
\[
f_S(\rho(x)) \;=\; r_S-\frac{\langle \vw_S,x\rangle}{\|\vw_S\|\|x\|},
\]
so the inequalities above becomes
\[
y\in \rho(S)\Rightarrow f_S(y)\le 0,\qquad y\in Y\setminus \rho(S)\Rightarrow f_S(y)>0.
\]
Therefore we show that
\[
\med(m, k; \fln) <= \med(m, k; \fcos)
\]

Conversely, let $n^\star=\mathrm{MED}(m,k;\fln)$. Then there exists $X=\{x_1,\dots,x_m\}\subset\R^{n^\star}$ such that for every $S\subseteq X$ with $|S|\le k$ there is $f_S(x)=\langle \vw_S,x\rangle+b_S$ satisfying
\[
x\in S \Rightarrow f_S(x)\le 0,\qquad x\in X\setminus S \Rightarrow f_S(x)>0.
\]

Consider the embedding $\phi:\R^{n^\star}\to S^{n^\star}\subset\R^{n^\star+1}$ given by
\[
\phi(x)\;=\;\frac{(x,1)}{\|(x,1)\|},
\]
where $(x,1)$ denotes the concatenation of $x$ and 1.

Let $Y=\{\phi(x_i)\}_{i=1}^m\subset S^{n^\star}$. For each $S$, set
\[
u_S\;=\;\frac{(\vw_S,b_S)}{\|(\vw_S,b_S)\|}\in S^{n^\star},\qquad t_S:=0,
\]
and define $g_S(y)=\tfrac{\langle u_S,y\rangle}{\|u_S\|\|y\|}-t_S=\langle u_S,y\rangle$ (since $\|u_S\|=\|y\|=1$ on $S^{n^\star}$).
Then for every $x\in\R^{n^\star}$,
\[
g_S(\phi(x))
=\left\langle \frac{(\vw_S,b_S)}{\|(\vw_S,b_S)\|},\,\frac{(x,1)}{\|(x,1)\|}\right\rangle
=\frac{\langle \vw_S,x\rangle+b_S}{\|(\vw_S,b_S)\|\cdot\|(x,1)\|}.
\]

Therefore, for $x\in X$ we have the exact shattering
\[
f_S(x)\le 0 \iff g_S(\phi(x))\le 0,\qquad
f_S(x)>0 \iff g_S(\phi(x))>0.
\]
\end{proof}

\subsection{Proof of Theorem~\ref{thm:medc-linear}, Theorem~\ref{thm:medc-cos}, and Theorem~\ref{thm:medc-l2}}\label{app:proof:thm:maed-linear}

Those three proofs are very similar, so they are grouped together.

\begin{proof}
Let $v_1, \dots, v_m \sim \mathcal{N}(0, I_n/n)$ in $\real^n$, our goals are to show that there is a positive probability such that
for any subset $S$, $|S|\leq k$, $v_1\in S$ and $v_2\notin S$, the following inequalities holds:
\begin{align}
    \langle v_1, \sum_{u\in S} u \rangle & > \langle v_2 ,  \sum_{u\in S} u \rangle, & \text{for Theorem~\ref{thm:medc-linear}};\\
    \langle \frac{v_1}{\|v_1\|}, \sum_{u\in S} u \rangle & > \langle \frac{v_2}{\|v_2\|},  \sum_{u\in S} u \rangle, &\text{for Theorem~\ref{thm:medc-cos}};\\
    \|v_1 - \frac{1}{|S|} \sum_{u\in S} u \|_2^2 &< \|v_2 - \frac{1}{|S|} \sum_{u\in S} u \|_2^2, & \text{for Theorem~\ref{thm:medc-l2}}.
\end{align}

To show this, considering the concentration of the inner product of two independent random vectors $v_i$ and $v_j$ drawn from $\mathcal{N}(0, \frac{1}{n})$:
\begin{align}
\Pr\left(|\langle v_i, v_j\rangle| \geq \frac{1}{3k} \right) \leq 2\exp\left(-c\frac{n}{k^2}\right),
\end{align}
where $c$ is a constant.

Also, noticing that the norm of such vectors is concentrated as
\begin{align}
\Pr\left(|\|v_i\| - 1| \geq \frac{1}{3k} \right) \leq 2\exp\left(-c\frac{n}{k^2}\right),
\end{align}

\textbf{To prove theorem~\ref{thm:medc-linear}}, let $k=l$.
For $m$ vectors, take the union bound of (1) inner products of all possible pairs and (2) the norm of all vectors, then we derive the probability that any of the inner products (or the $\|v_i - 1\|$) is greater than $\frac{1}{3k}$, and then force it to be smaller than 1.
\begin{align}
	2 \left(m+\frac{m(m-1)}{2}\right) \exp\left(-c\frac{n}{k^2}\right) < 1.
\end{align}

The probability of the event where any inner products are smaller than $\frac{1}{3k}$, and norms are no diverge from 1 larger than $\frac{1}{3k}$ is positive when
\begin{align}
	n > C k^2 \log m.
\end{align}
where $C$ is a constant.

Under such conditions,
\begin{align}
    &LHS = \langle v_1, \sum_{u\in S} u \rangle = \|v_1\| + \sum_{u\in S-\{v_1\}} \langle v_1, u\rangle > 1 - \frac{1}{3k} + (|S|-1) \frac{1}{3k} = 1 - \frac{|S|}{3k}, \\
    &RHS = \sum_{u\in S} \langle v_2, u\rangle < \frac{|S|}{3k}.
\end{align}
Because $|S| \leq k$, so $LHS > RHS$, Theorem~\ref{thm:medc-linear} is proved.

\textbf{To prove Theorem~\ref{thm:medc-cos}}
\begin{align}
    LHS &= \langle \frac{v_1}{\|v_1\|}, \sum_{u\in S} u \rangle > \frac{1 - |S| / 3k}{1 + 1/3k} = \frac{3k-|S|}{3k + 1} \\
    RHS &= \langle \frac{v_2}{\|v_2\|}, \sum_{u\in S} u \rangle < \frac{|S|}{3k(1-\frac{1}{3k})} = \frac{|S|}{3k-1}.
\end{align}
Then,
\begin{align}
    LHS - RHS  > \frac{3k(3k - 2|S| -1)}{9k^2-1}.
\end{align}
Because $1 < |S| \leq k$, so $LHS - RHS > 0$ and Theorem~\ref{thm:medc-cos} is proved.

\textbf{To prove Theorem~\ref{thm:medc-l2}}
\begin{align}
    LHS & = \|v_1 - \frac{1}{|S|} \sum_{u\in S} u \|_2^2 = \|v_1\|_2^2 + \|\frac{1}{|S|} \sum_{u\in S} u \|_2^2 - \frac{2}{|S|}\langle v_1, \sum_{u\in S} u \rangle\\
    & < 1+\frac{1}{3k} - \frac{2}{|S|}(1-\frac{|S|}{3k})  + \|\frac{1}{|S|} \sum_{u\in S} u \|_2^2, \\
    RHS & = \|v_2 - \frac{1}{|S|} \sum_{u\in S} u \|_2^2 = \|v_2\|_2^2 + \|\frac{1}{|S|} \sum_{u\in S} u \|_2^2 - \frac{2}{|S|}\langle v_2, \sum_{u\in S} u \rangle \\
    & >  1-\frac{1}{3k} - \frac{2}{|S|}\frac{|S|}{3k} + \|\frac{1}{|S|} \sum_{u\in S} u \|_2^2.
\end{align}
Then,
\begin{align}
    LHS - RHS < \frac{2}{3k} + \frac{2}{3k} - \frac{2}{|S|}(1-\frac{|S|}{3k}) = \frac{2}{k} - \frac{2}{|S|}.
\end{align}
Because $|S| \leq k$, so $LHS - RHS < 0$ and Theorem~\ref{thm:medc-l2} is proved.
\end{proof}
